\documentclass[conference]{IEEEtran}
\usepackage{amsmath,booktabs,array}
\usepackage{url}
\title{Runtime Hardware Trojan Detection and Regional Localization in an INT8 FPGA-CNN Accelerator}
\author{\IEEEauthorblockN{Jeevan K. S.}\IEEEauthorblockA{Department of Electronics and Communication Engineering\\IIIT Dharwad, India}}
\begin{document}\maketitle
\begin{abstract}
This paper presents a lightweight runtime Hardware Trojan detection and regional localization framework for an INT8 LeNet-style CNN accelerator implemented on an Intel Cyclone V FPGA. Five controlled Trojan variants target Conv2 processing-element arithmetic, weight memory, feature interconnect, spatial routing, and control sequencing. A canonical 60-row cycle-accurate simulation matrix covers ten MNIST workloads and Healthy/T1--T5 targets. A separate held-out same-CNN evaluation selected 50 Healthy-correct images (five per class) and ran each on all six targets, yielding 300 additional simulation rows. The canonical matrix gives TP=50, TN=10, FP=0 and FN=0; the held-out matrix gives TP=250, TN=50, FP=0 and FN=0. Both results are scoped to their controlled, Healthy-correct workload sets. Regional localization is correct for all 50 Trojan rows at the defined three-region resolution. T1--T4 detection latency is 149,136 cycles (2.98272 ms at 50 MHz), while T5 is detected in 149,135 cycles. Fresh C6 Quartus evidence is retained for Healthy, T1 and T2. T1--T5 were physically programmed on a Terasic DE10-Standard with the expected class/detector/localization LED pattern. Quartus power values are tool estimates, not physical rail measurements.
\end{abstract}
\begin{IEEEkeywords}
hardware Trojan, FPGA security, CNN accelerator, runtime detection, regional localization, INT8, MNIST
\end{IEEEkeywords}

\section{Introduction}
Hardware Trojans are malicious hardware modifications that may remain dormant or preserve externally visible behavior on selected tests. Classical work identifies the difficulty of detecting small or rarely activated modifications using conventional verification alone \cite{chakraborty2009hardware}. Neural-network accelerators introduce additional attack surfaces because repeated MACs, memories, interconnects and control paths can be modified while the final class may remain unchanged for selected inputs \cite{clements2018hardware}. A survey of DNN hardware security similarly identifies Trojan, side-channel and fault-injection threats as important security dimensions \cite{mittal2021survey}.

This work demonstrates an end-to-end runtime defense for an FPGA CNN. The contributions are: (1) an INT8 LeNet-style reference accelerator; (2) five controlled Trojan implementations; (3) a compact detector with regional localization; (4) a deterministic ten-workload simulation matrix; and (5) Quartus and physical-board evidence with explicit scientific scope boundaries.

\section{Related Work}
Prior work relevant to this study falls into four groups. First, neural-network and CNN accelerator Trojan attacks show that internal feature maps, arithmetic and reconfigurable interconnects can be manipulated while preserving selected output behavior. Clements and Lao introduced neural-network Trojan attacks, while FeSHI studied feature-map-level stealthy attacks; more recent work considers reconfigurable FPGA-CNN interconnections and PUF-based countermeasures \cite{clements2018hardware,odetola2021feshi,hou2024interconnect}. A two-level protection scheme and hardware-resilient CNN designs pursue mitigation through architectural protection or redundancy rather than the regional runtime monitor studied here \cite{liu2024twolevel,sun2022resilient,sun2025resilient}.

Second, runtime and anomaly-based approaches monitor circuit behavior or functional relationships. Changepoint-based detection and reverse-function runtime monitoring provide relevant precedent for runtime detection, but they do not evaluate the same five-region INT8 FPGA-CNN campaign used here \cite{elnaggar2019changepoint,mohd2021rmvrf}. Third, structural localization methods use netlist/RTL graph information or gate-level learning. Golden-reference-free graph-convolutional localization and TrojanSAINT illustrate structural approaches; explainable LUT-level localization advances fine-grained diagnosis but does not provide the same on-board runtime regional code \cite{yasaei2022localization,lashen2023trojansa,su2025lutlocalization}. Fourth, lightweight secure accelerator architectures, including Nacc-Guard, seek to reduce attack impact through design-level defenses \cite{li2023naccguard}.

The distinction claimed here is not that these methods are inferior. Rather, this work targets an end-to-end evidence layer combining a small runtime detector, three-region localization, cycle-accurate workload regression, Quartus evidence, and physical DE10-Standard output validation. Direct numerical ranking is inappropriate because the cited studies use different datasets, threat models, attack payloads, hardware platforms, and metrics.

\section{Architecture and Threat Model}
The reference network accepts 28$\times$28 grayscale MNIST images and contains Conv1, ReLU, Pool1, Conv2, ReLU, Pool2, a fully connected classifier and ArgMax. The computation is
\begin{equation}
Y=\operatorname{ArgMax}(f_{FC}(f_{Pool2}(f_{Conv2}(f_{Pool1}(f_{Conv1}(X)))))).
\end{equation}
The attacker is assumed to have RTL/netlist-level access during an untrusted design or synthesis stage. The controlled campaign includes arithmetic corruption (T1), weight-memory corruption (T2), feature-interconnect alteration (T3), source-routing alteration (T4), and a one-cycle control stall (T5).

The detector output is mapped to LEDR[3]. A two-bit localization code is mapped to LEDR[5:4]: 00=healthy, 01=Conv2 PE arithmetic, 10=Conv2 weight memory, and 11=shared interconnect/routing/control. LEDR[2:0] reports the predicted class. Code 11 intentionally does not distinguish T3, T4 and T5.

\section{Experimental Methodology}
Cycle-accurate Verilog simulation uses Icarus Verilog. The canonical matrix contains ten workloads (digits 0--9) and six targets, giving 60 rows. The FPGA target is the Intel Cyclone V SoC 5CSXFC6D6F31C6 on the Terasic DE10-Standard. The verified synchronous-M10K reference latency is 634,281 cycles at 50 MHz. Earlier 301,854-cycle and 346,563-cycle values are historical development milestones and are not used as the current reference.

\section{Results}
The canonical simulation ledger contains 10 Healthy and 50 Trojan rows. Therefore TP=50, TN=10, FP=0 and FN=0:
\begin{equation}
TPR=\frac{TP}{TP+FN},\quad FPR=\frac{FP}{FP+TN},
\end{equation}
\begin{equation}
Precision=\frac{TP}{TP+FP},\quad F1=2\frac{Precision\cdot TPR}{Precision+TPR}.
\end{equation}
For the canonical matrix, $(TP,TN,FP,FN)=(50,10,0,0)$, so the metrics evaluate to $(TPR,FPR,Precision,F1)=(1,0,1,1)$. The held-out matrix independently gives $(250,50,0,0)$ and the same point estimates within its Healthy-screened sample.
These metrics apply only to the canonical committed controlled matrix. In a separate held-out same-CNN evaluation, 50 additional Healthy-correct images were run on Healthy and T1--T5 (300 rows), producing TP=250, TN=50, FP=0 and FN=0. Because images misclassified by the Healthy baseline were screened out, the held-out FPR is conditional on the selected Healthy-correct set. The successful CI artifact and screening protocol are documented in docs/EXTENDED_HELDOUT_RESULTS.md.

T1 is localized to region 01 for 10/10 workloads, T2 to region 10 for 10/10, and T3--T5 to shared region 11 for 30/30. Regional localization accuracy is therefore 50/50=100\% at the defined resolution.

At 50 MHz, T1--T4 are detected after 149,136 cycles = 2.98272 ms. T5 is detected after 149,135 cycles = 2.98270 ms. T1--T4 inference is 634,281 cycles; T5 is 634,282 cycles.

\begin{table}[t]
\centering
\caption{Fresh C6 resource evidence}
\begin{tabular}{lrrrrrr}
\toprule
Target & ALM & Reg. & RAM & DSP & Setup & Hold\\
\midrule
Healthy & 1153 & 912 & 71 & 13 & +0.181 & +0.173\\
T1 & 1170 & 918 & 71 & 13 & +0.141 & +0.141\\
T2 & 1192 & 919 & 71 & 15 & +0.163 & +0.163\\
\bottomrule
\end{tabular}
\end{table}
Exact final fitter counts for T3--T5 and exact fresh-C6 Fmax are not archived and are not inferred.

Quartus Power Analyzer estimates are 463.79, 465.51, 466.28, 467.17, 464.75 and 464.27 mW for Healthy and T1--T5 respectively. These are tool estimates rather than physical rail/current measurements.

T1, T2, T3, T4 and T5 were programmed on the physical DE10-Standard with zero programming errors and zero warnings. The observed LED patterns matched the expected outputs: class 7 and detector asserted for all variants, localization 01 for T1, 10 for T2, and shared code 11 for T3--T5.

\section{Discussion and Limitations}
The experiment shows that the selected controlled Trojans can be detected and regionally localized even when the evaluated class output remains correct. The 100\% result is not a universal Trojan-detection guarantee. Limitations include absence of physical power instrumentation, incomplete archived final fitter counts for T3--T5, lack of archived fresh-C6 Fmax, a single compact MNIST CNN, and an evaluation-layer rather than independently synthesized numerical ablation.

\section{Conclusion}
A reproducible runtime detection and regional localization prototype was implemented for an INT8 FPGA CNN. The controlled simulation campaign achieved 100\% detection and regional localization at the defined scope, while physical DE10 validation was completed for T1--T5. The repository separates simulation metrics, implementation evidence, tool-estimated power and physical observations to avoid overstating the results.

\bibliographystyle{IEEEtran}
\bibliography{references/references}
\end{document}
